\section{Lines in the Cartesian Plane}

A line should not first appear as a formula to be recognised. We already have
all the ingredients needed to construct it geometrically: a point tells us where
to begin, and a nonzero vector tells us in which direction motion is allowed.
The question is therefore:

\begin{center}
\emph{Which points can be reached from one fixed point by moving only in one
fixed direction?}
\end{center}

\begin{corebox}[Data that determine a line]
A line in the Cartesian plane is determined by one point $P_0$ on the line and
one nonzero direction vector $\mathbf v$. Every point of the line is reached by
adding a scalar multiple of $\mathbf v$ to $P_0$.
\end{corebox}

Let
\[
P_0=(x_0,y_0)
\]
be fixed and let
\[
\mathbf v=(u,v)\neq\mathbf0
\]
be a fixed direction vector. Every permitted displacement is a scalar multiple
$t\mathbf v$. Starting from $P_0$ gives
\[
P=P_0+t\mathbf v,
\qquad t\in\mathbb R.
\]

\begin{definitionbox}
The line through $P_0$ with direction vector $\mathbf v\neq\mathbf0$ is
\[
\boxed{
L=\{P_0+t\mathbf v:t\in\mathbb R\}.
}
\]
In coordinates,
\[
\boxed{
\begin{aligned}
x&=x_0+tu,\\
y&=y_0+tv,
\end{aligned}
\qquad t\in\mathbb R.
}
\]
\end{definitionbox}

\begin{interpretationbox}[Meaning of the parameter]
At $t=0$ we are at $P_0$; positive and negative values move in the two opposite
orientations along the same direction. Changing the length of $\mathbf v$ only
changes how quickly the parameter runs through the points, not the geometric
line itself.
\end{interpretationbox}

\begin{center}
\begin{tikzpicture}[scale=0.9,>=stealth]
    \draw[axis] (-0.8,0) -- (6.2,0) node[right] {$x$};
    \draw[axis] (0,-0.8) -- (0,4.8) node[above] {$y$};
    \coordinate (P) at (1.2,1.1);
    \coordinate (Q) at (4.8,3.5);
    \coordinate (D) at (3.0,2.3);
    \draw[subset] (0.3,0.5) -- (5.6,4.03) node[above right] {$L$};
    \fill[geometricSubsetColor] (P) circle (2.2pt)
        node[below right] {$P_0$};
    \draw[->,blue!70!black,line width=1.1pt] (P) -- (D)
        node[midway,above left] {$\mathbf v$};
    \fill[geometricSubsetColor] (Q) circle (2.2pt)
        node[above left] {$P_0+t\mathbf v$};
\end{tikzpicture}
\end{center}

\subsection{The same line described by a perpendicular direction}

A direction vector tells us how to move along the line. There is a second,
equally important description: choose a nonzero vector $\mathbf n=(a,b)$
perpendicular to $\mathbf v$. A point $P=(x,y)$ lies on the line precisely when
the displacement $P-P_0$ is perpendicular to $\mathbf n$.

\begin{derivationbox}[From direction to normal form]
The dot-product test from the previous chapter gives
\[
\mathbf n\cdot(P-P_0)=0.
\]
Written in coordinates,
\[
a(x-x_0)+b(y-y_0)=0.
\]
After collecting the constant terms,
\[
\boxed{ax+by=c},
\qquad c=ax_0+by_0.
\]
\end{derivationbox}

\begin{interpretationbox}
The general linear equation is not a second, unrelated definition of a line. It
records the same object by specifying a direction that no displacement along the
line may contain.
\end{interpretationbox}

\begin{center}
\begin{tikzpicture}[scale=0.9,>=stealth]
    \coordinate (P0) at (1.1,1.1);
    \coordinate (P) at (4.6,3.2);
    \coordinate (N) at (0.0,3.1);
    \draw[subset] (0.2,0.56) -- (5.5,3.74) node[above right] {$L$};
    \draw[blue!70!black,line width=1.1pt,->] (P0) -- (P)
        node[midway,below right] {$P-P_0$};
    \draw[orange!85!black,line width=1.1pt,->] (P0) -- (N)
        node[midway,left] {$\mathbf n$};
    \draw (P0) ++(0.28,0.17) -- ++(-0.16,0.27) -- ++(-0.28,-0.17);
    \fill[geometricSubsetColor] (P0) circle (2.2pt) node[below left] {$P_0$};
    \fill[geometricSubsetColor] (P) circle (2.2pt) node[above] {$P$};
\end{tikzpicture}
\end{center}

\subsection{Slope form as a special coordinate description}

\begin{derivationbox}[Eliminating the parameter]
If the horizontal component $u$ of the direction vector is nonzero, eliminate
$t$ from
\[
x=x_0+tu,
\qquad
y=y_0+tv.
\]
Since
\[
t=\frac{x-x_0}{u},
\]
we obtain
\[
y-y_0=\frac vu(x-x_0).
\]
Therefore the slope is
\[
\boxed{m=\frac vu},
\]
and after collecting constants the equation takes the familiar form
\[
y=mx+b.
\]
\end{derivationbox}

\begin{remarkbox}
If $u=0$, the line is vertical and is described by $x=x_0$. The vector and
normal descriptions treat this case without any exception; only slope notation
fails.
\end{remarkbox}

\begin{summarybox}
The three forms answer different questions about the same set:
\[
\begin{array}{c|c}
P_0+t\mathbf v & \text{how to move along the line},\\
\mathbf n\cdot(P-P_0)=0 & \text{which displacements are allowed},\\
y=mx+b & \text{how vertical position depends on horizontal position}.
\end{array}
\]
No representation is the line itself. The line is the geometric set of points
that all these descriptions identify.
\end{summarybox}
